\chapter{Package and Pinout}
\label{ch:package}

The sound device is a single integrated circuit with a memory of its own attached. The sound processor, the registers
and the slot scheduler are on the one piece of silicon. The audio memory is a separate static memory that the device
alone controls, and it is not on the CPU's bus. The signals that leave the package are therefore the bus, one interrupt
output, the clock and reset inputs, the memory interface, and the serial output for the converter. This chapter gives
the groups of signals and the meaning of each. Chapter~\ref{ch:signal-timing} gives the timing of the signals.

\section{Package}

The package and the pin numbers are not yet specified. The device has 77 signal pins, listed in
Table~\ref{tab:signals}, in addition to its power and ground pins, whose number is also not yet specified.

\section{Signal Groups}

Figure~\ref{fig:signal-groups} shows the signals by group. Three groups face three neighbours: the CPU's bus, the audio
memory, and the converter.

\begin{figure}[H]
	\centering
	\begin{tikzpicture}[
			scale=0.9, transform shape,
			group/.style={draw, rounded corners=2pt, fill=sirclightgray, align=center, font=\small, minimum width=3.3cm, minimum height=1.5cm},
			chip/.style={draw, thick, fill=gray!8, align=center, font=\bfseries, minimum width=3.0cm, minimum height=5.6cm},
			flow/.style={-{Stealth[length=2mm]}, thick},
			note/.style={font=\scriptsize, text=sircgray, align=center}
		]
		\node[chip] (chip) at (0,0) {SIRCAS};
		\node[group] (bus) at (-6.2,1.9) {Bus interface\\A0--A15, D0--D15\\\activelow{CS}, \activelow{BAS}, BRW,\\\activelow{BACK}};
		\node[group] (clk) at (-6.2,-0.2) {Clock and reset\\CLKI, \activelow{RST}};
		\node[group] (irq) at (-6.2,-2.1) {Interrupt\\\activelow{IRQ3}};
		\node[group] (mem) at (6.2,1.6) {Audio memory interface\\MA0--MA14, MD0--MD15\\\activelow{MCE}, \activelow{MOE}, \activelow{MWE}};
		\node[group] (dac) at (6.2,-1.9) {Converter interface\\BCLK, WS, SD, SCK};
		\draw[flow, <->] (bus.east) -- (bus.east -| chip.west);
		\draw[flow, <-] (clk.east) -- (clk.east -| chip.west);
		\draw[flow] (chip.west |- irq) -- (irq.east);
		\draw[flow, <->] (chip.east |- mem) -- (mem.west);
		\draw[flow] (chip.east |- dac) -- (dac.west);
	\end{tikzpicture}
	\caption{SIRCAS Signal Groups}
	\label{fig:signal-groups}
\end{figure}

\section{Signal Descriptions}

Table~\ref{tab:signals} lists the direction and sense of each signal. Active-low signals carry a bar over the name;
they are asserted when low. Address, data and converter signals are active high. An output marked \emph{open drain} can
only pull its line low; the system supplies a pull-up, and several devices may share the line.

\begin{table}[H]
	\centering
	\small
	\begin{tabularx}{\textwidth}{llllX}
		\toprule
		\textbf{Signal}  & \textbf{Count} & \textbf{Direction} & \textbf{Output type} & \textbf{Function}                                              \\
		\midrule
		CLKI             & 1              & Input              &                      & Rising edge advances the device one cycle                      \\
		\activelow{RST}  & 1              & Input              &                      & Resets the device                                              \\
		\activelow{CS}   & 1              & Input              &                      & Selects the device for a bus cycle                             \\
		\activelow{BAS}  & 1              & Input              &                      & Marks a bus cycle                                              \\
		BRW              & 1              & Input              &                      & High for a read, low for a write                               \\
		\activelow{BACK} & 1              & Output             & Open drain           & Completes the bus cycle                                        \\
		A0--A15          & 16             & Input              &                      & Word offset in the device segment                              \\
		D0--D15          & 16             & Bidirectional      & Three-state          & Data                                                           \\
		\activelow{IRQ3} & 1              & Output             & Open drain           & A completed sequence is waiting (Section~\ref{sec:interrupts}) \\
		MA0--MA14        & 15             & Output             &                      & Word address of the audio memory                               \\
		MD0--MD15        & 16             & Bidirectional      & Three-state          & Audio memory data                                              \\
		\activelow{MCE}  & 1              & Output             &                      & Audio memory chip enable                                       \\
		\activelow{MOE}  & 1              & Output             &                      & Audio memory output enable                                     \\
		\activelow{MWE}  & 1              & Output             &                      & Audio memory write enable                                      \\
		BCLK             & 1              & Output             &                      & Serial bit clock to the converter                              \\
		WS               & 1              & Output             &                      & Serial word select: low for left, high for right               \\
		SD               & 1              & Output             &                      & Serial data                                                    \\
		SCK              & 1              & Output             &                      & Converter system clock                                         \\
		\bottomrule
	\end{tabularx}
	\caption{SIRCAS Signals}
	\label{tab:signals}
\end{table}

\subsection{Clock Input (CLKI)}

\textbf{Direction:} Input. The device is clocked by \texttt{CLKI}, a continuous clock of 10,738,636 Hz in an NTSC
system: half the frequency of the 21,477,272 Hz system oscillator. Every interval in this manual that is measured in
cycles is a count of \texttt{CLKI} cycles. The sample period is 336 cycles (Chapter~\ref{ch:timing}). The device
works at any steady frequency from 9.0 to 12.0~MHz (Section~\ref{sec:clock-rate}). In a PAL system \texttt{CLKI} is 10,640,685~Hz, half the 21,281,370~Hz system oscillator.

\subsection{Reset Input (RST)}

\textbf{Direction:} Input. \textbf{Signal Type:} Active Low. While \activelow{RST} is asserted
the device holds every register in its reset state (Section~\ref{sec:reset}). The audio memory is not cleared: its
contents are undefined after power-up and after reset.

\subsection{Bus Signals}

\activelow{BAS}, BRW, A0--A15, D0--D15 and \activelow{BACK} are bus signals of the SIRC-1, defined in the \hwmanual{}. \activelow{CS} is produced by the system's address decoding of the device's segment. A cycle is addressed to the device when \activelow{CS} and \activelow{BAS} are both
asserted. The address is a word offset in the device's 65,536-word segment (Chapter~\ref{ch:memory-map}). The device
asserts \activelow{BACK} to complete the cycle. Register accesses are completed without delay. An access to the audio
memory window is completed when the device has made it, which takes between two and five \texttt{CLKI} cycles
(Section~\ref{sec:cpu-slots}).

\subsection{Interrupt Output (IRQ3)}
\label{sec:pin-irq3}

\textbf{Direction:} Output. \textbf{Signal Type:} Active Low, open drain. The device holds \activelow{IRQ3} asserted
for as long as at least one sequence has completed whose interrupt is enabled (Section~\ref{sec:interrupts}). The
signal is level-sensitive. It is released when no enabled flag remains set.

\subsection{Audio Memory Interface (MA, MD, MCE, MOE, MWE)}
\label{sec:memory-interface}

The audio memory is 32,768 words of 16 bits, so it has a 15-bit address and a 16-bit data bus. It may be built from two
32,768 by 8 bit memories used side by side. The memory interface belongs to the sound device alone: the CPU's bus and
the memory's bus are separate, and the device is the only user of the memory. The device makes every access, both its
own and the CPU's (Section~\ref{sec:cpu-slots}), in a fixed pattern of slots, so the CPU is never in competition with
the sound processor for the memory.

Every memory transaction, a read or a write, occupies two \texttt{CLKI} cycles, which is 186 ns. The memory must
therefore complete an access within two cycles. The electrical figures of the interface are in
Section~\ref{sec:electrical}; the timing of \activelow{MCE}, \activelow{MOE} and \activelow{MWE} within an access is
not yet specified.

\subsection{Converter Interface (BCLK, WS, SD, SCK)}
\label{sec:converter-interface}

The device is the master of the serial interface to the converter. It sends the two stereo channels time-multiplexed on
\texttt{SD}, with \texttt{WS} low for the left channel and high for the right, and \texttt{BCLK} clocking the bits
(Chapter~\ref{ch:signal-timing}). \texttt{SCK} is a clock for converters that need a system clock in addition to the
bit clock; its frequency and phase relation to \texttt{BCLK} are not yet specified.

\section{Electrical Characteristics}
\label{sec:electrical}

The figures in this section are provisional. They are targets the device is designed to meet, and will be confirmed or
revised when it is characterised. They are the same as those of the CPU (\hwmanual{}, electrical characteristics
chapter), so that the three devices fit together on one bus. The supply is 5~V $\pm$ 5\%, the signals are TTL
compatible, and all figures apply over 0 to 70~\textdegree{}C. Times are measured from or to the 1.4~V point of the
rising edge of \texttt{CLKI}, with the outputs loaded by 50~pF.

\begin{table}[H]
	\centering
	\begin{tabular}{p{7.6cm}lrrl}
		\toprule
		\textbf{Parameter}                                                                         & \textbf{Symbol}    & \textbf{Min} & \textbf{Max} & \textbf{Unit} \\
		\midrule
		\texttt{CLKI} frequency                                                                    & $f_{CLKI}$         & 9.0          & 12.0         & MHz           \\
		\texttt{CLKI} high time and low time                                                       & $t_{CH}$, $t_{CL}$ & 35           &              & ns            \\
		Input high voltage, input low voltage                                                      &                    & 2.0, 0       & VCC, 0.8     & V             \\
		Output high voltage ($I_{OH} = -400~\mu$A), output low voltage ($I_{OL} = 4$~mA)           &                    & 2.4          & 0.4          & V             \\
		Low level of an open-drain output, $I_{OL} = 8$~mA                                         &                    &              & 0.4          & V             \\
		Input capacitance                                                                          &                    &              & 10           & pF            \\
		\midrule
		Bus inputs (A, D on a write, BRW, \activelow{CS}, \activelow{BAS}): set-up before the edge & $t_{IS}$           & 10           &              & ns            \\
		Bus inputs: hold after the edge                                                            & $t_{IH}$           & 0            &              & ns            \\
		Read data and \activelow{BACK} valid after \activelow{CS} and \activelow{BAS} asserted     & $t_{AV}$           &              & 35           & ns            \\
		Read data float, after the edge that ends a read                                           & $t_{OZ}$           &              & 25           & ns            \\
		Output hold                                                                                & $t_{OH}$           & 3            &              & ns            \\
		\bottomrule
	\end{tabular}
	\caption{AC and DC Characteristics}
	\label{tab:electrical}
\end{table}

\noindent The memory interface outputs (MA, MCE, MOE, MWE and MD on a write) are valid within 25~ns of the edge that starts the access, and are held for 3~ns, and the read data has 10~ns of set-up before the edge that ends it, with no hold. An access takes two cycles, 186~ns, so the memory's address access time may be at most 140~ns, with 5~ns left for the board. At 12~MHz an access is 167~ns and the limit is 120~ns.
